% Einheit 1 – Terme aufstellen und berechnen (bank/terme/e1.jsonl, zusammenbau v1.1)
\einheitenkopf*[e1][Terme aufstellen und berechnen]{Terme aufstellen und berechnen}
\setcounter{aufgabe}{20}


% test: terme-e1-k1-s7-v2
\lbtest{Kannst du das schon? Dann bist du fertig}{
\lbtt{a)}{Die Summe aus dem Vierfachen einer Zahl $x$ und 9 wird halbiert. Kreuze den passenden Term an. \\ $\square$\ $4x + 9 : 2$ \\ $\square$\ $(4x + 9) : 2$ \\ $\square$\ $2 \cdot (4x + 9)$ \\ $\square$\ $4x : 9 + 2$}
}{a) $(4x + 9) : 2$}


% ohne: terme-e1-k2-s1-v1, terme-e1-k2-s1-v2, terme-e1-k2-s1-v3
\begin{kbaufgabe}{Termwerte berechnen}
\begin{kbblock}
\kbteil{a)}{$4x - 3$ für $x = 5$: \lbfeld}{}
\end{kbblock}
\begin{kbblock}
\kbteil{b)}{$2x + 9$ für $x = -3$: \lbfeld}{}
\end{kbblock}
\begin{kbblock}
\kbteil{c)}{$10 - 3x$ für $x = -2$: \lbfeld}{}
\end{kbblock}
\end{kbaufgabe}


% leiter: terme-e1-k1-s1-v1, terme-e1-k1-s1-v2, terme-e1-k1-s2-v1, terme-e1-k1-s3-v1, terme-e1-k1-s5-v1, terme-e1-k1-s6-v1
\begin{kbaufgabe}{Einen Term aufstellen}
\begin{kbblock}
\kbanweisung{Kreuze den Term an, der dazu passt.}
\kbteil{a)}{Das Vierfache einer Zahl $x$ wird um 3 vermindert.}{}
\kbkreuz{$4x - 3$}
\kbkreuz{$4 \cdot (x - 3)$}
\kbkreuz{$3 - 4x$}
\end{kbblock}
\begin{kbblock}
\kbteil{b)}{Das Vierfache einer Zahl $x$ wird um 5 vermindert.}{}
\kbkreuz{$5 - 4x$}
\kbkreuz{$4x - 5$}
\kbkreuz{$4 \cdot (x - 5)$}
\end{kbblock}
\begin{kbblock}
\kbteil{c)}{Das Dreifache einer Zahl $x$ wird um 8 vergrößert. Schreibe dazu einen Term.}{}
\kbantwort{\kbfeld}
\end{kbblock}
\begin{kbblock}
\kbteil{d)}{Die Figur zeigt ein Rechteck. Die Maße sind in cm angegeben. Schreibe einen Term für den Flächeninhalt des Rechtecks.}{}
\kbzweispaltig[0.42]{\viereck[seiten={5,x,5,x}]{(0,0)}{(5,0)}{(5,2.5)}{(0,2.5)}}{\lbantwortrechts[1.5cm]{\kbfeld}}
\end{kbblock}
\begin{kbblock}
\kbteil{e)}{Ein Rechteck hat die Seiten $x$ cm und 9 cm. Schreibe einen Term für den Umfang. Fasse den Term zusammen.}{}
\kbantwort{\kbfeld}
\end{kbblock}
\begin{kbblock}
\kbteil{f)}{Eine Rechnung geht so: Verdreifache eine Zahl $x$, addiere 2 und verdopple die Summe. Schreibe dazu einen Term.}{}
\kbantwort{\kbfeld}
\end{kbblock}
\end{kbaufgabe}


% pruefung: terme-e1-k1-s7-v1, terme-e1-k1-s7-v3
\begin{kbaufgabe}{Wie in der Prüfung}
\begin{kbblock}
\kbteil{a)}{Die Differenz aus dem Dreifachen einer Zahl $x$ und 7 wird verdoppelt. Kreuze den passenden Term an.}{P10 ’23}
\kbkreuz{$(3x - 7) : 2$}
\kbkreuz{$3x : 7 \cdot 2$}
\kbkreuz{$2 \cdot (3x - 7)$}
\kbkreuz{$3x - 7 \cdot 2$}
\item[]\kbraum{2}
\end{kbblock}
\begin{kbblock}
\kbteil{b)}{Das Sechsfache einer Zahl $x$ wird um zwei vermindert. Kreuze den passenden Term an.}{P10 ’17}
\kbkreuz{$6 \cdot (x - 2)$}
\kbkreuz{$2 - 6x$}
\kbkreuz{$6x - 2$}
\kbkreuz{$6x - 2x$}
\item[]\kbraum{2}
\end{kbblock}
\end{kbaufgabe}


% ohne: terme-e1-k3-s1-v1, terme-e1-k3-s1-v2, terme-e1-k3-s1-v3
\begin{kbaufgabe}{Zu einem Term eine passende Situation erfinden}
\begin{kbblock}
\kbanweisung{Erfinde eine Situation aus dem Alltag, zu der der Term … passt. Schreibe auf, wofür $x$ steht. Berechne den Term für … und erkläre, was das Ergebnis bedeutet.}
\kbteil{a)}{$15 + 3x$, $x = 2$}{}
\item[]\kbraum{2}
\end{kbblock}
\begin{kbblock}
\kbteil{b)}{$20 - 2x$, $x = 4$}{}
\item[]\kbraum{2}
\end{kbblock}
\begin{kbblock}
\kbteil{c)}{$10 + 4x$, $x = 3$}{}
\item[]\kbraum{2}
\end{kbblock}
\end{kbaufgabe}


% pflicht: terme-e1-k4-s1-v3, terme-e1-k4-s2-v2, terme-e1-k4-s3-v1, terme-e1-k4-s4-v1
\begin{lbkasten}
\begin{kbaufgabe}{Verstanden?}
\begin{kbblock}
\item[]\lbhinweis{Gemischt – erkenne selbst, was zu tun ist.}
\kbteil{a)}{Vier Schüler haben einen Text in einen Term übersetzt. Genau eine der vier Zeilen ist falsch. Welche? Schreibe den Term richtig.}{}
\kbfrage{(1) das Vierfache einer Zahl $x$, um 5 vergrößert: $4x + 5$ \\ (2) das Doppelte der Summe aus $x$ und 6: $2x + 6$ \\ (3) eine Zahl $x$, vermindert um 8: $x - 8$ \\ (4) die Hälfte einer Zahl $x$: $x : 2$}
\item[]\kbraum{2}
\end{kbblock}
\begin{kbblock}
\kbteil{b)}{Entscheide bei jeder Aussage, ob sie wahr oder falsch ist. Begründe, ohne genau zu rechnen.}{}
\kbfrage{(1) „Eine Zahl $x$, vermindert um 3“ und „3 vermindert um eine Zahl $x$“ ergeben immer denselben Term. \\ (2) Das Doppelte einer Summe braucht immer eine Klammer. \\ (3) Es gibt eine Zahl $x$, für die $2 \cdot (x + 4)$ und $2x + 4$ denselben Wert haben.}
\item[]\kbraum{3}
\end{kbblock}
\begin{kbblock}
\kbteil{c)}{Welcher Term gehört zum Termbaum?}{}
\kbzweispaltig[0.38]{\termbaum{\cdot}{4}{\tb{+}{x}{6}}}{\lbantwortrechts[1.1cm]{\kbfeld}}
\end{kbblock}
\begin{kbblock}
\kbteil{d)}{Ein Stromtarif kostet 11 € Grundpreis im Monat und 0,30 € je Kilowattstunde. Stelle einen Term für die Monatskosten in Euro bei $n$ Kilowattstunden auf. Was kostet ein Monat mit 150 Kilowattstunden?}{}
\item[]\kbraum{3}
\end{kbblock}
\end{kbaufgabe}
\end{lbkasten}
